\documentclass[10pt,a4paper]{amsart}
\usepackage[margin=19mm]{geometry}
\usepackage{amsmath,amssymb,mathtools}
\usepackage[scr=rsfs,cal=euler]{mathalfa}
\usepackage{xfrac}
\usepackage[unicode,hidelinks,bookmarks=false]{hyperref}
\newcommand{\ie}{i.e.\ }
\newcommand{\cf}{cf.\ }
\newcommand{\resp}{respectively\ }
\newcommand{\Subsection}{\subsection}
\providecommand{\nobreakdash}{\nobreak\hbox{-}}
\setlength{\emergencystretch}{2em}
\multlinegap0pt
\usepackage{mathtools}
\usepackage{amssymb}
\newtheorem{lemma}{Lemma}
\newcommand{\hE}{e^{\mathrm{uc}}_{\mathrm{top}}}
\newcommand{\huc}{h^{\mathrm{uc}}_{\mathrm{top}}}
\title{Entropy on regular sets and periodic-orbit growth for singular flows}
\author{C.A. Morales}
\date{Source: arXiv:2609.00626v1 (CC BY 4.0)}
\begin{document}
\maketitle

\noindent\emph{Source and licence.} Entropy on regular sets and periodic-orbit growth for singular flows. Version arXiv:2609.00626v1 (CC BY 4.0), distributed by arXiv under the Creative Commons Attribution 4.0 International licence, http://creativecommons.org/licenses/by/4.0/, as recorded in the arXiv licence metadata for this exact version and shown on the abstract page https://arxiv.org/abs/2609.00626v1. Full licence notice: \texttt{licenses/arxiv-2609.00626-CC-BY-4.0.txt}.\par\smallskip

\noindent\textit{Editorial scope.} Counted unit: the article's lemma l1 (source0.tex lines 427--435). The article's complete original proof of it is delivered with it (source0.tex lines 437--566, one \texttt{proof} environment of 130 lines). No mathematics is written, repaired or re-derived: the statement and the proof are the article's own lines, in the article's own notation, and no retained line is substituted or re-flowed.  No further source-local context is needed: every cross-reference of every retained line resolves inside the two delivered spans.  Nothing was omitted from any retained span, so the package carries no omission record.  No retained line contains a citation, so the wrapper carries no bibliography and no bibliography entry.  No retained line contains a figure environment, an image-inclusion command or drawing code, so no image is delivered and the package declares no asset dependency.  Editorial formatting only, in addition: an AMS \texttt{amsart} preamble that restates the article's own declarations and macros used by the retained lines, namely the environments \texttt{lemma} on separate counters, together with the standard packages those lines need.  The retained environments print in this document as Lemma 1 in order of appearance, whereas the article prints the same items with the numbering of its own full document.  The complete native source of the article as distributed in the arXiv e-print for this exact version is bundled unchanged as \texttt{sources/209045/source0.tex}.
\begin{lemma}
\label{l1}
For every compact set $K\subset X$,
\[
\hE(\varphi,K)
=
\huc(f,K).
\]
\end{lemma}

\begin{proof}
We give the proof using spanning sets.
Since the integer times
\[
0,1,\ldots,n-1
\]
belong to the interval $[0,n]$, one has
\[
d_n^f(x,y)\leq d_n^\varphi(x,y).
\]
Therefore every $(n,\varepsilon)$-spanning set for the flow is also
$(n,\varepsilon)$-spanning for the time-one map. Hence
\[
r_n^f(K,\varepsilon)
\leq
r_n^\varphi(K,\varepsilon).
\]
It follows that
\[
\lim_{\varepsilon\to0}
\limsup_{n\to\infty}
\frac1n\log r_n^f(K,\varepsilon)
\leq
\huc(\varphi,K).
\]

We now prove the reverse inequality. Fix $\varepsilon>0$. Since the
map
\[
[0,1]\times X\longrightarrow X,
\qquad
(t,x)\longmapsto\varphi_t(x),
\]
is continuous on a compact space, it is uniformly continuous.
Therefore there exists $\delta=\delta(\varepsilon)>0$ such that
\[
d(x,y)<\delta
\quad\Longrightarrow\quad
d\bigl(\varphi_u(x),\varphi_u(y)\bigr)<\varepsilon
\]
for every $u\in[0,1]$. We may choose
\[
0<\delta(\varepsilon)\leq\varepsilon,
\]
so that $\delta(\varepsilon)\to0$ as $\varepsilon\to0$.

Suppose that
\[
d_{n+1}^f(x,y)<\delta.
\]
Let $t\in[0,n]$. Write
\[
t=j+u,
\]
where
\[
j\in\{0,\ldots,n-1\},
\qquad
u\in[0,1].
\]
Then
\[
\varphi_t(x)=\varphi_u(f^j(x))
\]
and
\[
\varphi_t(y)=\varphi_u(f^j(y)).
\]
Since
\[
d\bigl(f^j(x),f^j(y)\bigr)<\delta,
\]
uniform continuity gives
\[
d\bigl(\varphi_t(x),\varphi_t(y)\bigr)<\varepsilon.
\]
Taking the maximum over $t\in[0,n]$, we obtain
\[
d_n^\varphi(x,y)<\varepsilon.
\]

Thus every $(n+1,\delta)$-spanning set for $K$ with respect to the
time-one map is $(n,\varepsilon)$-spanning for $K$ with respect to the
flow. Consequently,
\[
r_n^\varphi(K,\varepsilon)
\leq
r_{n+1}^f(K,\delta).
\]
Therefore
\[
\begin{aligned}
\limsup_{n\to\infty}
\frac1n\log r_n^\varphi(K,\varepsilon)
&\leq
\limsup_{n\to\infty}
\frac1n\log r_{n+1}^f(K,\delta)\\
&=
\limsup_{n\to\infty}
\frac1n\log r_n^f(K,\delta).
\end{aligned}
\]
Letting $\varepsilon\to0$, and hence
$\delta(\varepsilon)\to0$, gives
\[
\hE(\varphi,K)
\leq
\lim_{\delta\to0}
\limsup_{n\to\infty}
\frac1n\log r_n^f(K,\delta).
\]
Together with the opposite inequality, this proves
\[
\hE(\varphi,K)
=
\lim_{\varepsilon\to0}
\limsup_{n\to\infty}
\frac1n\log r_n^f(K,\varepsilon).
\]

Finally, the standard inequalities
\[
r_n^f(K,\varepsilon)
\leq
s_n^f(K,\varepsilon)
\leq
r_n^f\left(K,\frac{\varepsilon}{2}\right)
\]
show that the separated-set formula gives the same entropy.
\end{proof}

\end{document}
